% GENERATED FILE. Edit canonical lessons, then run npm run generate:books.
% canonical-source-hash: fnv1a64:464db525bee97bde
% canonical-lessons: HI-C531-dholak, HI-C531-sahnai, HI-C531-ghughru, HI-C531-baja, HI-C531-mrdang

\chapter{Drums and Pipes}
\label{ch:hi-drums-and-pipes}
\hlchaptermodality{531}

\begin{chapteropening}
\textbf{By the end of this chapter:} \emph{I can name instruments and say who plays them.}

Complete the last lesson of chapter 531: I can name instruments and say who plays them.
\end{chapteropening}

\section[\texorpdfstring{\hi{ढोलक} (ḍholak)}{ढोलक (ḍholak)}]{\hi{ढोलक} (ḍholak) — a small two-headed drum}

\label{lesson:HI-C531-dholak}



\begin{quote}
Before the new one: say the Hindi for a raga, then the Hindi for to hum.
\end{quote}

\subsection*{You'll want to know: \hi{ढोलक}}
\textbf{\hi{ढोलक}} — \emph{ḍholak} — \textquotedblleft{}a small two-headed drum\textquotedblright{}.

\begin{quote}
\hi{शादी} \hi{में} \hi{औरतें} \hi{ढोलक} \hi{बजाकर} \hi{गाती} \hi{हैं।}
\end{quote}

\begin{quote}
\emph{Śādī meṁ aurteṁ ḍholak bajākar gātī haiṁ.}
\end{quote}

\begin{quote}
At weddings the women play the dholak and sing.
\end{quote}

Feminine.

\begin{cousinweb}
\textbf{\hi{ढोलक}} (\emph{ḍholak}) is a tadbhava, inherited from Sanskrit \emph{ḍhōla}, \textquotedblleft{}a drum\textquotedblright{} (CDIAL 5608).
\end{cousinweb}

\begin{grammarlens}[title={What you've built}]
One more word for describing what you lived and felt.
\end{grammarlens}

\subsection*{Your turn}
\begin{itemize}
\raggedright
  \item \emph{Say it:} \emph{ḍholak}
  \item \emph{Say it:} \emph{ḍholak}, once more
  \item \emph{Say it:} \emph{Śādī meṁ aurteṁ ḍholak bajākar gātī haiṁ.} — At weddings the women play the dholak and sing.
  \item \emph{Recall:} say the Hindi for a raga, then the Hindi for to hum, then say \emph{ḍholak} again
\end{itemize}

\begin{tcolorbox}[breakable,colback=teal!4,colframe=teal!35!black,title={Before you move on}]
What does \hi{ढोलक} mean? (\textquotedblleft{}A small two-headed drum\textquotedblright{}.) Say it once more.
\end{tcolorbox}
\section[\texorpdfstring{\hi{शहनाई} (śahnāī)}{शहनाई (śahnāī)}]{\hi{शहनाई} (śahnāī) — a shehnai (a reed pipe)}

\label{lesson:HI-C531-sahnai}



\begin{quote}
Before the new one: say the Hindi for to hum, then the Hindi for a small two-headed drum.
\end{quote}

\subsection*{You'll want to know: \hi{शहनाई}}
\textbf{\hi{शहनाई}} — \emph{śahnāī} — \textquotedblleft{}a shehnai (a reed pipe)\textquotedblright{}.

\begin{quote}
\hi{शादी} \hi{के} \hi{दिन} \hi{सुबह} \hi{से} \hi{शहनाई} \hi{बज} \hi{रही} \hi{थी।}
\end{quote}

\begin{quote}
\emph{Śādī ke din subah se śahnāī baj rahī thī.}
\end{quote}

\begin{quote}
On the wedding day the shehnai had been playing since morning.
\end{quote}

Feminine.

\begin{cousinweb}
\textbf{\hi{शहनाई}} (\emph{śahnāī}) is usually traced to Persian \emph{shahnāy}, \textquotedblleft{}a king's pipe\textquotedblright{}.
\end{cousinweb}

\begin{grammarlens}[title={What you've built}]
One more word for describing what you lived and felt.
\end{grammarlens}

\subsection*{Your turn}
\begin{itemize}
\raggedright
  \item \emph{Say it:} \emph{śahnāī}
  \item \emph{Say it:} \emph{śahnāī}, once more
  \item \emph{Say it:} \emph{Śādī ke din subah se śahnāī baj rahī thī.} — On the wedding day the shehnai had been playing since morning.
  \item \emph{Recall:} say the Hindi for to hum, then the Hindi for a small two-headed drum, then say \emph{śahnāī} again
\end{itemize}

\begin{tcolorbox}[breakable,colback=teal!4,colframe=teal!35!black,title={Before you move on}]
What does \hi{शहनाई} mean? (\textquotedblleft{}A shehnai (a reed pipe)\textquotedblright{}.) Say it once more.
\end{tcolorbox}
\section[\texorpdfstring{\hi{घुँघरू} (ghũghrū)}{घुँघरू (ghũghrū)}]{\hi{घुँघरू} (ghũghrū) — ankle bells}

\label{lesson:HI-C531-ghughru}



\begin{quote}
Before the new one: say the Hindi for a small two-headed drum, then the Hindi for a shehnai (a reed pipe).
\end{quote}

\subsection*{You'll want to know: \hi{घुँघरू}}
\textbf{\hi{घुँघरू}} — \emph{ghũghrū} — \textquotedblleft{}ankle bells\textquotedblright{}.

\begin{quote}
\hi{उसके} \hi{पैरों} \hi{के} \hi{घुँघरू} \hi{बज} \hi{रहे} \hi{थे।}
\end{quote}

\begin{quote}
\emph{Uske pairoṁ ke ghũghrū baj rahe the.}
\end{quote}

\begin{quote}
The bells on her feet were ringing.
\end{quote}

Masculine.

\begin{cousinweb}
\textbf{\hi{घुँघरू}} (\emph{ghũghrū}) is a tadbhava, inherited from \emph{*ghuṅghura}, \textquotedblleft{}a string of bells\textquotedblright{}, a form CDIAL reconstructs from its descendants (CDIAL 4477).
\end{cousinweb}

\begin{grammarlens}[title={What you've built}]
One more word for describing what you lived and felt.
\end{grammarlens}

\subsection*{Your turn}
\begin{itemize}
\raggedright
  \item \emph{Say it:} \emph{ghũghrū}
  \item \emph{Say it:} \emph{ghũghrū}, once more
  \item \emph{Say it:} \emph{Uske pairoṁ ke ghũghrū baj rahe the.} — The bells on her feet were ringing.
  \item \emph{Recall:} say the Hindi for a small two-headed drum, then the Hindi for a shehnai (a reed pipe), then say \emph{ghũghrū} again
\end{itemize}

\begin{tcolorbox}[breakable,colback=teal!4,colframe=teal!35!black,title={Before you move on}]
What does \hi{घुँघरू} mean? (\textquotedblleft{}Ankle bells\textquotedblright{}.) Say it once more.
\end{tcolorbox}
\section[\texorpdfstring{\hi{बाजा} (bājā)}{बाजा (bājā)}]{\hi{बाजा} (bājā) — a musical instrument; a band}

\label{lesson:HI-C531-baja}



\begin{quote}
Before the new one: say the Hindi for a shehnai (a reed pipe), then the Hindi for ankle bells.
\end{quote}

\subsection*{You'll want to know: \hi{बाजा}}
\textbf{\hi{बाजा}} — \emph{bājā} — \textquotedblleft{}a musical instrument; a band\textquotedblright{}.

\begin{quote}
\hi{मेले} \hi{में} \hi{बाजा} \hi{बज} \hi{रहा} \hi{था।}
\end{quote}

\begin{quote}
\emph{Mele meṁ bājā baj rahā thā.}
\end{quote}

\begin{quote}
A band was playing at the fair.
\end{quote}

Masculine.

\begin{cousinweb}
\textbf{\hi{बाजा}} (\emph{bājā}) is a tadbhava, inherited from Sanskrit \emph{vādya}, \textquotedblleft{}to be played; music\textquotedblright{} (CDIAL 11511).
\end{cousinweb}

\begin{grammarlens}[title={What you've built}]
One more word for describing what you lived and felt.
\end{grammarlens}

\subsection*{Your turn}
\begin{itemize}
\raggedright
  \item \emph{Say it:} \emph{bājā}
  \item \emph{Say it:} \emph{bājā}, once more
  \item \emph{Say it:} \emph{Mele meṁ bājā baj rahā thā.} — A band was playing at the fair.
  \item \emph{Recall:} say the Hindi for a shehnai (a reed pipe), then the Hindi for ankle bells, then say \emph{bājā} again
\end{itemize}

\begin{tcolorbox}[breakable,colback=teal!4,colframe=teal!35!black,title={Before you move on}]
What does \hi{बाजा} mean? (\textquotedblleft{}A musical instrument; a band\textquotedblright{}.) Say it once more.
\end{tcolorbox}
\section[\texorpdfstring{\hi{मृदंग} (mṛdaṅg)}{मृदंग (mṛdaṅg)}]{\hi{मृदंग} (mṛdaṅg) — a mridang (a barrel drum)}

\label{lesson:HI-C531-mrdang}



\begin{quote}
Before the new one: say the Hindi for ankle bells, then the Hindi for a musical instrument.
\end{quote}

\subsection*{You'll want to know: \hi{मृदंग}}
\textbf{\hi{मृदंग}} — \emph{mṛdaṅg} — \textquotedblleft{}a mridang (a barrel drum)\textquotedblright{}.

\begin{quote}
\hi{मंदिर} \hi{में} \hi{एक} \hi{बूढ़ा} \hi{आदमी} \hi{मृदंग} \hi{बजा} \hi{रहा} \hi{था।}
\end{quote}

\begin{quote}
\emph{Mandir meṁ ek būṛhā ādmī mṛdaṅg bajā rahā thā.}
\end{quote}

\begin{quote}
In the temple an old man was playing the mridang.
\end{quote}

Masculine.

\begin{cousinweb}
\textbf{\hi{मृदंग}} (\emph{mṛdaṅg}) is a tatsama, taken whole from Sanskrit \emph{mṛdaṃga}, \textquotedblleft{}a kind of drum\textquotedblright{} (Monier-Williams).
\end{cousinweb}

\begin{grammarlens}[title={What you've built}]
One more word for describing what you lived and felt.
\end{grammarlens}

\subsection*{Your turn}
\begin{itemize}
\raggedright
  \item \emph{Say it:} \emph{mṛdaṅg}
  \item \emph{Say it:} \emph{mṛdaṅg}, once more
  \item \emph{Say it:} \emph{Mandir meṁ ek būṛhā ādmī mṛdaṅg bajā rahā thā.} — In the temple an old man was playing the mridang.
  \item \emph{Recall:} say the Hindi for ankle bells, then the Hindi for a musical instrument, then say \emph{mṛdaṅg} again
\end{itemize}

\begin{tcolorbox}[breakable,colback=teal!4,colframe=teal!35!black,title={Before you move on}]
What does \hi{मृदंग} mean? (\textquotedblleft{}A mridang (a barrel drum)\textquotedblright{}.) Say it once more.
\end{tcolorbox}
